\section{Résultats}\label{sec:resultats}

\subsection{Comportement du mur}
Le profil de température à travers le mur de trois couches (figure~\ref{fig:mur}) montre que l'isolant
concentre la chute de température : entre l'interface béton/isolant et l'interface isolant/plâtre, la
température passe de 34,36 à 25,58~\textdegree C, soit environ 8,8~\textdegree C sur les 10~\textdegree C
d'écart entre l'extérieur et l'intérieur.

\begin{figure}[ht]
	\centering
	\includegraphics[width=\textwidth]{cas2_profil}
	\caption{Mur de trois couches en régime stationnaire : profil de température à $y=0{,}5$~m (courbe : calcul
		P1 ; points : solution exacte) et carte de température.}
	\label{fig:mur}
\end{figure}

En régime transitoire (figure~\ref{fig:transitoire}), après un échelon de la température extérieure à
35~\textdegree C, la face extérieure atteint 33,7~\textdegree C au bout de 12~h, alors que la face intérieure
ne passe que de 25 à 25,26~\textdegree C : le béton absorbe la chaleur et l'isolant protège l'intérieur. Après
96~h, le régime stationnaire est atteint.

\begin{figure}[ht]
	\centering
	\includegraphics[width=\textwidth]{cas3_temps}
	\caption{Régime transitoire du mur après un échelon de température extérieure : température en quatre
		points (les lignes pointillées sont les valeurs du régime stationnaire) et flux thermique.}
	\label{fig:transitoire}
\end{figure}

Sur une journée type (figure~\ref{fig:cyclemur}), le mur amortit 98,5~\% de l'oscillation extérieure
(amplitude de 0,075~\textdegree C à la face intérieure contre 5~\textdegree C dehors) et retarde son maximum
de 7,7~h : la chaleur de l'après-midi atteint la face intérieure vers 22~h~40.

\begin{figure}[ht]
	\centering
	\includegraphics[width=\textwidth]{cas5_journee}
	\caption{Cycle jour/nuit sur le mur (régime périodique) : température en quatre points sur 24~h et flux
		thermique entrant et sortant.}
	\label{fig:cyclemur}
\end{figure}

\subsection{Pièce en 2D : ponts thermiques}
La carte de température de la pièce de référence (figure~\ref{fig:piece}) fait apparaître l'effet des
jonctions. Loin des coins, la surface intérieure est à 25,44~\textdegree C, comme dans le calcul 1D du mur ;
aux coins elle est plus chaude : 26,64~\textdegree C aux coins supérieurs (jonction mur/toit) et
28,06~\textdegree C aux coins inférieurs (jonction mur/plancher). Le milieu du plancher reste à
25,0~\textdegree C.

Le flux total entrant vaut 46,6~W/m, contre 33,7~W/m prédits par un calcul 1D sur les 9,6~m de murs et de toit
intérieurs ($3{,}51~\mathrm{W/m^2}\times 9{,}6~\mathrm{m}$). Les jonctions et le plancher ajoutent donc environ
12,9~W/m, soit 28~\% du flux total. Ce supplément s'explique vraisemblablement par la continuité du béton
entre la dalle, non isolée, et la couche extérieure des murs.

\begin{figure}[ht]
	\centering
	\includegraphics[width=\textwidth]{cas4_carte}
	\caption{Pièce de référence en régime stationnaire : carte de température et zooms sur les coins du côté
		gauche (jonction mur/toit en haut, jonction mur/plancher en bas).}
	\label{fig:piece}
\end{figure}

\subsection{Scénarios en régime stationnaire}
Cinq scénarios sont comparés sur la pièce, avec un intérieur maintenu à 25~\textdegree C : la référence (S0),
la suppression de l'isolant (S1, couche remplacée par du béton), un isolant plus performant
($\lambda=0{,}025$~W/(m\,K), S2), un toit au soleil (température équivalente de 54,2~\textdegree C, S3) et un
toit clair (44,6~\textdegree C, S4). Le tableau~\ref{tab:scen} donne les résultats.

\begin{table}[ht]
	\centering\small
	\begin{tabular}{lcccc}
		\toprule
		Scénario & flux total (W/m) & écart & T moy. surfaces (\textdegree C) & T max surfaces (\textdegree C) \\
		\midrule
		S0 référence & 46,6 & -- & 25,43 & 28,06 \\
		S1 sans isolant & 243,6 & $+423$~\% & 27,24 & 29,50 \\
		S2 isolant performant & 35,1 & $-25$~\% & 25,32 & 28,07 \\
		S3 toit au soleil & 74,0 & $+59$~\% & 25,68 & 28,06 \\
		S4 toit clair & 60,3 & $+29$~\% & 25,55 & 28,06 \\
		\bottomrule
	\end{tabular}
	\caption{Scénarios en régime stationnaire : flux entrant total par mètre de profondeur et températures des
		surfaces intérieures.}
	\label{tab:scen}
\end{table}

Dans chaque scénario, le flux calculé au milieu d'un mur et au milieu du toit est égal à la valeur 1D
(3,51 ; 23,85 ; 2,30 ; 10,25 et 6,88~W/m$^2$ selon le cas). En retirant du flux total la part des murs
et du toit, on obtient la contribution des jonctions et du plancher (tableau~\ref{tab:decomp},
figure~\ref{fig:flux}).

\begin{table}[ht]
	\centering\small
	\begin{tabular}{lcccc}
		\toprule
		Scénario & murs (W/m) & toit (W/m) & jonctions et plancher (W/m) & total (W/m) \\
		\midrule
		S0 & 19,66 & 14,05 & 12,89 & 46,60 \\
		S1 & 133,57 & 95,40 & 14,60 & 243,56 \\
		S2 & 12,88 & 9,20 & 13,02 & 35,10 \\
		S3 & 19,66 & 41,01 & 13,36 & 74,04 \\
		S4 & 19,66 & 27,53 & 13,13 & 60,32 \\
		\bottomrule
	\end{tabular}
	\caption{Décomposition du flux total entrant (murs : $2\times2{,}8$~m, toit : $4{,}0$~m).}
	\label{tab:decomp}
\end{table}

\begin{figure}[ht]
	\centering
	\includegraphics[width=0.9\textwidth]{scenarios_flux}
	\caption{Flux entrant total par scénario, décomposé en murs, toit, jonctions et plancher.}
	\label{fig:flux}
\end{figure}

\begin{itemize}
	\item Supprimer l'isolant multiplie par 6,8 le flux à travers un mur (de 3,51 à 23,85~W/m$^2$).
	\item Un isolant plus performant réduit le flux à travers un mur de 35~\%, mais le flux total de 25~\%
	seulement : la contribution des jonctions, d'environ 13~W/m (de 12,9 à 14,6~W/m selon le scénario),
	ne dépend presque pas de l'isolant des murs. La température maximale des surfaces intérieures reste
	de ce fait voisine de 28~\textdegree C au coin inférieur, sauf sans isolant (29,5~\textdegree C).
	\item Le soleil sur le toit ajoute 59~\% au flux total. Une peinture claire divise exactement cette hausse
	par deux ($+13{,}7$~W/m au lieu de $+27{,}4$~W/m).
\end{itemize}

\begin{figure}[ht]
	\centering
	\includegraphics[width=\textwidth]{scenarios_cartes}
	\caption{Cartes de température des cinq scénarios, avec la même échelle de couleurs.}
	\label{fig:cartes}
\end{figure}

\subsection{Confort sur une journée}
Dans les simulations de cycle jour/nuit, l'intérieur n'est plus climatisé : l'air évolue librement et la
surchauffe est mesurée par les degrés-heures $DH_{28}$ et $DH_{30}$ de l'équation~\eqref{eq:dh}. Cinq scénarios
sont comparés : la référence (C0, absorption du toit de 0,6 et un renouvellement d'air par heure), la
suppression de l'isolant (C1), un isolant plus performant (C2), un toit clair (C3, absorption de 0,3) et une
ventilation nocturne (C4, cinq renouvellements d'air par heure de 20~h à 6~h). Le tableau~\ref{tab:confort} et
la figure~\ref{fig:confort} donnent les résultats sur le dernier jour simulé.

\begin{table}[ht]
	\centering\small
	\begin{tabular}{lccccccc}
		\toprule
		Scénario & T air (\textdegree C) & écart (\textdegree C) & ampl. (\textdegree C) & max (h) & $DH_{28}$ & $DH_{30}$ & écart $DH_{30}$ \\
		\midrule
		C0 référence & 31,04 & -- & 0,55 & 18,83 & 72,9 & 24,9 & -- \\
		C1 sans isolant & 32,13 & $+1{,}09$ & 0,77 & 21,33 & 99,1 & 51,1 & $+105$~\% \\
		C2 isolant performant & 30,79 & $-0{,}25$ & 0,51 & 18,50 & 66,9 & 18,9 & $-24$~\% \\
		C3 toit clair & 30,52 & $-0{,}52$ & 0,50 & 18,67 & 60,5 & 12,5 & $-50$~\% \\
		C4 ventilation nocturne & 29,41 & $-1{,}63$ & 1,25 & 20,00 & 33,8 & 2,4 & $-90$~\% \\
		\bottomrule
	\end{tabular}
	\caption{Confort sur 24 h (régime périodique) : température moyenne de l'air, écart à la référence, amplitude,
		heure du maximum et degrés-heures par jour ($^\circ$C\,h).}
	\label{tab:confort}
\end{table}

\begin{figure}[ht]
	\centering
	\includegraphics[width=\textwidth]{cycle_comparaison}
	\caption{Température de l'air intérieur sur 24 h pour les cinq scénarios (à gauche) et degrés-heures
		au-dessus de 30~\textdegree C (à droite).}
	\label{fig:confort}
\end{figure}

\begin{itemize}
	\item Dans tous les scénarios, l'air oscille très peu (amplitude de 0,5 à 1,25~\textdegree C contre
	5~\textdegree C à l'extérieur) : la masse du bâtiment lisse la journée, et le maximum de l'air
	intervient en fin d'après-midi ou en soirée.
	\item De C0 à C3, l'air reste au-dessus de 28~\textdegree C toute la journée : $DH_{28}$ vaut alors
	simplement $24\times(\overline{T}_{\mathrm{air}}-28)$ et ne mesure que le niveau moyen. Le seuil de
	30~\textdegree C est plus discriminant.
	\item Le toit clair divise exactement par deux l'excès de chaleur au-dessus de 30~\textdegree C. Ce
	résultat s'explique par le modèle : sans apports internes, l'air serait à la moyenne extérieure
	(30~\textdegree C), et tout l'excès provient du rayonnement absorbé par le toit, proportionnel à
	l'absorption (0,6 contre 0,3).
	\item La ventilation nocturne est la mesure la plus efficace : $-1{,}63$~\textdegree C en moyenne,
	$-90$~\% de degrés-heures, avec un air qui descend jusqu'à 28,1~\textdegree C. Elle exploite des
	nuits plus fraîches (jusqu'à 25~\textdegree C) pour rafraîchir l'air et la masse du bâtiment.
	\item Sans isolant, l'air est plus chaud de 1,1~\textdegree C, les degrés-heures doublent et le maximum
	est décalé à 21~h~20.
\end{itemize}

Ces classements sont valables dans le cadre des hypothèses du modèle : le classement des solutions est plus
fiable que les valeurs absolues. Un point mérite d'être signalé pour la ventilation nocturne : à 20~h,
l'extérieur (31,3~\textdegree C) est encore plus chaud que l'air intérieur (environ 30,6~\textdegree C), si
bien que la ventilation forte fait d'abord entrer de la chaleur, ce qui explique l'amplitude plus grande
(1,25~\textdegree C) et le maximum à 20~h. Une ventilation commandée, déclenchée seulement quand l'extérieur
est plus frais que l'intérieur, ferait mieux ; les petites marches visibles sur la courbe C4 à 6~h et à 20~h
viennent du changement brutal de débit combiné à un air sans inertie propre.